\documentclass[12pt]{article}
\usepackage[utf8]{inputenc}
\usepackage[T1]{fontenc}
\usepackage{hyperref}
\title{Kyanga Language Resources for K-12 Curriculum Development}
\author{Compiled by Claude Sonnet 5 (claude-sonnet-5) for Andr\'as Kornai}
\date{}
\begin{document}
\maketitle

\section{Introduction}
Kyanga (Kyenga, Tyenga; ISO 639-3 tye) is an Eastern Mande language of Nigeria and Benin. Sources disagree on its vitality and speaker numbers. Joshua Project gives 28,000 speakers in Nigeria and 2,500 in Benin. Wikipedia gives about 12,000.

\section{Primary Educational Resources}
No primers or school materials were found.

\section{Linguistic and Reference Materials}
\begin{itemize}
\item Wikipedia, Kyenga language: \url{https://en.wikipedia.org/wiki/Kyenga_language}
\item DoBeS (MPI Nijmegen) Kyanga and Shanga documentation project pages: \url{https://dobes.mpi.nl/projects/kyanga_shanga/geography/}, \url{https://dobes.mpi.nl/?p=3054}, \url{https://dobes.mpi.nl/?p=3052}
\item DoBeS poster on Kyanga: \url{https://dobes.mpi.nl/posters/print/africa-kyanga-print.pdf}
\end{itemize}

\section{Religious and Community Resources}
Scripture Earth page for tye: \url{https://www.scriptureearth.org/00eng.php?iso=tye}. The search did not show whether any scripture is actually available. Joshua Project: \url{https://www.joshuaproject.net/index.php/languages/tye}.

\section{Digital and Audio Resources}
Wikitongues page: \url{https://wikitongues.org/languages/tye/}. The DoBeS project is a documentation project; the search did not show what recordings it holds.

\section{Curriculum Development Context}
No Kyanga school materials were found. Hausa in Nigeria and Dendi in Benin are displacing the language. The DoBeS documentation is the most promising source for further material.
\end{document}